%% kapa: exam=1 round=35 year=2024 period=1 form=A subject=real_estate subject_ko=부동산학원론 seq=5 number=85
%% source: Q-Net 감정평가사 기출문제 (cst00309) artlSeq=5213229 file=raw/1008/5213229/1교시 문제지 원본.pdf
%% conversion: kordoc 4.24.51 -> tools/gen1.py (auto); review=proofread-sonnet; flags=none
\begin{question}{85}
건축원자재 가격의 하락에 따른 영향을 디파스퀠리-위튼(DiPasquale \& Wheaton)의 사분면 모형을 통해 설명한 것으로 옳지 않은 것은? (단, 주어진 조건에 한함)
\begin{choices}
  \item 건축원자재 가격의 하락으로 인해 부동산개발부문에서 신규건설비용이 하락한다.
  \item 주어진 부동산자산가격 수준에서 부동산개발의 수익성이 높아지므로 신규건설량이 증가한다.
  \item 새로운 장기균형에서 균형공간재고는 감소한다.
  \item 새로운 장기균형에서 부동산공간시장의 균형임대료는 하락한다.
  \item 새로운 장기균형에서 부동산자산시장의 균형가격은 하락한다.
\end{choices}
\end{question}
\answer{3}
